\section{The Preparatory Cycle: Presence, Adoration, and Offering}

The received Fátima narrative begins before May 1917. These earlier reports illuminate the message's grammar, but their source status must remain plain. Lúcia described three mist-like or human-shaped appearances in 1915 in her Fourth Memoir of 1941. She described three 1916 apparitions of an angel most fully in memoirs written in 1937 and 1941. No comparable contemporaneous 1915--1916 dossier has been located, and the 1930 operative sentence names the children's Marian visions from May through October 1917, not the preparatory cycle.

\subsection{The 1915 occurrences}

Lúcia recalled that, while praying with three girls on the hill, she saw on several occasions a white, nearly transparent human-like form above the trees. No words, identity, or mission were reported. In the later narrative these occurrences function as a threshold: perception before interpretation, an indistinct sign before a named messenger. Their very vagueness argues against treating every remembered phenomenon as a doctrinal communication.

\subsection{The three reported encounters with the Angel}

\begin{fourstable}{Encounter}{Narrative and action}{Spiritual content}{Documentary boundary}
Spring 1916, Loca do Cabeço & A luminous youth identifies himself as the Angel of Peace, prostrates, and leads the children in a threefold prayer. & Faith, adoration, hope, love, and intercession for those who lack them; the human person placed bodily before God. & Full text comes from the later memoir, not a 1916 written deposition.\\
Summer 1916, Poço do Arneiro & The angel interrupts play, calls for much prayer, and teaches the children to make a sacrifice of whatever can be offered. & Mercy, conversion of sinners, peace for Portugal, and patient acceptance of unavoidable suffering. & “Offer everything” must be read through Catholic moral prudence, never as permission for self-harm.\\
Autumn 1916, Loca do Cabeço & A host and chalice are presented in a Eucharistic vision; the angel adores the Trinity, gives Communion, and calls for reparation. & Real presence, Christ's Body and Blood, adoration, sin's ecclesial wound, conversion, and the Hearts of Jesus and Mary. & A visionary Communion narrative does not create a new sacramental rite or authorize reception apart from Church discipline.\\
\end{fourstable}

The first prayer can be summarized by its four verbs: \emph{believe, adore, hope, love}. They name the theological virtues with adoration at their center. The children then ask pardon for those who do not live these acts. Intercession does not divide humanity into the pure and the contemptible. The petitioner identifies with a wounded world and brings unbelief, indifference, despair, and lovelessness before mercy.

The second encounter changes “sacrifice” from a rare heroic event into a form of attention. Work, delay, hunger, insult, disappointment, and unavoidable pain can be offered rather than wasted. In Catholic doctrine, the offering has no independent redemptive value. Christ's sacrifice is unique, complete, and sufficient. Baptized persons can unite their lives to him because grace incorporates them into his Body (Rom 12:1; Col 1:24), not because Calvary lacks merit.

The third encounter is the cycle's summit. Its Trinitarian prayer adores Father, Son, and Holy Spirit and offers to the Father the Eucharistic Christ---Body, Blood, Soul, and Divinity---in reparation and supplication. The language is devotional rather than a second consecratory formula. Objectively, Christ offers himself once for all and sacramentally makes that sacrifice present in the Eucharist; the faithful join his self-offering. Subjectively, reparation means allowing his charity to oppose the disorder of sin through worship, repentance, and love.

\subsection{What “reparation” can and cannot mean}

Reparation does not appease a volatile deity by transferring pain to children. Sin really disorders the creature, wounds communion, profanes worship, and injures neighbors. God's love is not repaired as though divine beatitude could be diminished. The sinner and the world damaged by sin require healing and restored justice. Christ alone accomplishes reconciliation in its source; believers participate by repentance, restitution, adoration, intercession, mercy, and the patient offering of suffering that cannot rightly be removed.

“Consoling God” is analogical. God does not depend on creatures for emotional completion. Yet Scripture speaks truly of divine delight, grief, jealousy, and compassion according to the covenantal relation. To console Christ is to return love for love, remain with him in Gethsemane, honor him where he is despised, and care for the members with whom he identifies. Francisco's later Eucharistic language is sound when understood as communion, not as the repair of a deficient God.

The Angel's title “of Portugal” likewise needs proportion. Scripture can speak of angels associated with peoples (Dan 10), and Catholic tradition invokes patrons of places. No nation thereby becomes elect in place of the Church, morally innocent, or guaranteed a political program. Intercession for one's country becomes Catholic when it opens into prayer for enemies and universal peace.

\subsection{A school for the later message}

The preparatory cycle supplies a vertical architecture:

\begin{connectionstable}
Trinity & Every Marian element remains inside the worship of Father, Son, and Holy Spirit; Mary is neither destination nor divine source.\\
Eucharist & The message's center is Christ sacramentally present, not fascination with apparitions. The shrine is rightly ordered when it leads to Mass, confession, adoration, and charity.\\
Theological virtues & Faith, hope, and love judge every interpretation. A theory that breeds contempt, despair, or obsession has missed the announced purpose.\\
Offering & Ordinary life can be joined to Christ for others. The agency is grace-enabled and ecclesial, not magical transfer or spiritual coercion.\\
Peace & Peace begins in reconciliation with God but extends to nations and victims of war; it is neither passivity before injustice nor partisan victory.\\
\end{connectionstable}

This architecture also supplies a safety rule. Children may be taught prayer, generosity, fasting proportionate to age and health, and acceptance of unavoidable hardship. They may not be encouraged to conceal injury, seek pain, neglect food or medicine, obey abusive commands, or treat suffering as intrinsically desirable. The later September correction about a penitential rope---not while sleeping---already shows the received narrative placing zeal under restraint. Catholic asceticism is governed by charity, obedience, duties of state, bodily stewardship, and competent care.
